\documentclass[11pt,a4paper]{article}
\usepackage[margin=2.5cm]{geometry}
\usepackage{amsmath,amssymb}
\usepackage{booktabs}
\usepackage{listings}
\usepackage[colorlinks=true,linkcolor=blue,citecolor=blue,urlcolor=blue]{hyperref}

\lstset{basicstyle=\ttfamily\small,columns=fullflexible,frame=single,
        language=Fortran,keepspaces=true}

\newcommand{\kb}{\bar k}
\newcommand{\Nch}{N_{\rm ch}}
\newcommand{\nc}{{\rm nc}}

\title{Channel decomposition of the NLO real-emission integrand in MadNkLO}
\author{MadNkLO internal note}
\date{15 November 2026}

\begin{document}
\maketitle

\begin{abstract}
In each sector of the local analytic sector subtraction, the subtracted real
integrand $RW-K$ is integrated with a single phase-space mapping, the one of
the sector. The counterterm $K$, however, is a sum of pieces evaluated on
different mapped Born configurations. We describe an optional multi-channel
integration, switched on by \texttt{channel\_decomposition} in the run card,
in which $K$ is split by momentum mapping and the real matrix element is
shared among the channels with fractions computed from the counterterm. Each
channel is integrated with its own mapping and its own VEGAS grid. The sum over
channels is the original integrand at every phase-space point, both for the
total and for the histogram fills. We describe the construction, the
treatment of Born-level cuts (fractions computed from the uncut counterterm),
the extension to initial-state partons, and the validation for $e^+e^-\to$
jets, $pp\to e^+e^-$ and $pp\to Zj$.
\end{abstract}

\section{Motivation}

In sector $(i,j)$ the NLO real contribution to an observable $F$ is
\begin{equation}
  \sigma_{ij}[F] = \int d\Phi_{n+1}\,
  \Big[ R\,W_{ij}\,F(k) - \sum_{\rm pieces\ p} K_p\,F\big(\kb^{(p)}\big)\Big],
  \label{eq:sector}
\end{equation}
where $W_{ij}$ is the sector function and each counterterm piece $K_p$ is
evaluated on its own mapped Born configuration $\kb^{(p)}$. The cuts on
$\kb^{(p)}$ are included in $K_p$. By default the whole sector is integrated
with the mapping of the sector, $(i_U,i_S,i_B)$, which flattens the
singularities of $R\,W_{ij}$ in the limits of the sector.

The pieces of $K$ built with other mappings depend on Born configurations
$\kb^{(p)}$ that are not the integration variables. Cuts and bin edges on
these configurations are then not aligned with the VEGAS grid, and pieces
that cancel against $R\,W_{ij}$ in the integral can fill different bins at
nearby points. The channel decomposition is a standard multi-channel approach to this: each
piece is integrated with the mapping that generates its own Born
configuration.

\section{Construction}

\subsection{Channels}
For sector $(i,j)$, the set $\mathcal A_{ij}$ of channels is the set of
distinct momentum mappings $\alpha=(i_U,i_S,i_B)$ used by the counterterm:
\begin{itemize}
\item the soft dipoles $(i,l,m)$, with $l>m$ coloured and different from $i$
      (one mapping per dipole);
\item for the dipole between the two incoming partons, both initial--initial
      mappings $(i,1,2)$ and $(i,2,1)$, one channel each
      (sect.~\ref{sec:isr});
\item the hard-collinear piece $C_{ij}(1-S_i)$, including its soft-collinear
      subtraction, on the mapping of the sector $(i_U,i_S,i_B)$.
\end{itemize}
Mappings that occur more than once are listed once. For example, the
hard-collinear mapping often coincides with one of the soft dipoles. We write
\begin{equation}
  K = \sum_{\alpha\in\mathcal A_{ij}} K_\alpha ,
\end{equation}
where $K_\alpha$ is the sum of the pieces with mapping $\alpha$, with their cuts.
$K^\nc_\alpha$ denotes the same pieces with the cuts on their mapped Born
removed, and $K^\nc=\sum_\alpha K^\nc_\alpha$.

\subsection{Channel integrands}
Channel $\alpha$ integrates
\begin{equation}
  I_\alpha = \rho_\alpha\,R\,W_{ij} - K_\alpha ,
  \qquad
  \sum_{\alpha} \rho_\alpha = 1
  \;\;\Longrightarrow\;\;
  \sum_\alpha I_\alpha = R\,W_{ij} - K
  \label{eq:Ialpha}
\end{equation}
at every phase-space point. The integration uses the forward mapping
$\alpha$: a Born configuration $\kb^{(\alpha)}$ and three radiation variables
generate $k$, with the Jacobian of mapping $\alpha$. The channel fractions are
\begin{equation}
  \rho_\alpha =
  \begin{cases}
    K^\nc_\alpha / K^\nc
      & \text{if } \big|R\,W_{ij} - K^\nc\big| < \delta\,\big|K^\nc\big| ,\\[4pt]
    1/\Nch & \text{otherwise},
  \end{cases}
  \label{eq:rho}
\end{equation}
where $\delta$ is \texttt{channel\_delta} (default $1$) and $\Nch$ is the number
of channels of the sector.

\paragraph{Why these fractions.}
Near any singular limit the uncut counterterm reproduces the singular
behaviour of the real, $R\,W_{ij}\to K^\nc$, whatever the cuts. The criterion is
then satisfied and
\begin{equation}
  I_\alpha = \frac{K^\nc_\alpha}{K^\nc}\big(R\,W_{ij}-K^\nc\big)
           + \big(K^\nc_\alpha - K_\alpha\big).
  \label{eq:Ilimit}
\end{equation}
The first term is the subtracted real times a fraction that stays finite in
the limit, so it is integrable channel by channel. The second term is nonzero only where a cut on
$\kb^{(\alpha)}$ fails. It is the same Born-cut mismatch that exists without
channels, now shared among the channels in the same way as the singular part
of the real, which is the property that matters (sect.~\ref{sec:cuts}).
Hence the singular part of $R\,W_{ij}$ is distributed among the channels
exactly like the counterterm pieces that cancel it, and each channel is
integrable on its own.

Away from the limits any partition of unity is allowed. The ratio
$K_\alpha/K$ is not used there because the pieces have different signs:
where they cancel, $K\to0$ and the ratio would be unbounded. The fractions
$\rho_\alpha$ in eq.~\eqref{eq:rho} can lie outside $[0,1]$; only their sum
matters.

\paragraph{Luminosity.}
With hadron beams, the luminosity of $R$ and of every piece is taken at the
momentum fractions of the real-emission point, which phase\_space\_npo sets
for each channel. The fractions therefore form a partition of unity in each
beam orientation separately.

\subsection{Histograms and scale variations}
In channel $\alpha$:
\begin{itemize}
\item $\rho_\alpha R\,W_{ij}$ fills the histograms at the real-emission point
      $k$;
\item each piece of $K_\alpha$ fills at its own Born configuration
      $\kb^{(\alpha)}$.
\end{itemize}
The scale-variation weights get the same factors: $\alpha_s$ at the
real-emission or Born scale, and the PDF ratios of the $\mu_F$ variations of
each luminosity. Summed over channels, the fills are those of
eq.~\eqref{eq:sector}, point by point.

\subsection{Integration}
For each sector:
\begin{enumerate}
\item The warm-ups of all channels run first, with $N_{\rm grid}/\Nch$ calls
      each, and each channel's VEGAS grid is saved.
\item The $N$ calls of the main run are shared among the channels in
      proportion to the warm-up errors, with a minimum of $N/(10\Nch)$ per
      channel.
\item The results, histograms and scale variations of the channels are
      summed. The channels are independent runs, so their errors are added
      in quadrature.
\end{enumerate}
The total number of calls per sector is the same as without channels. The only
extra cost is that $K$ is evaluated twice (sect.~\ref{sec:cuts}).

\section{Born-level cuts: fractions from the uncut counterterm}
\label{sec:cuts}

The first implementation used $K_\alpha/K$, with cuts, both in the
fractions and in the criterion. In $pp\to Zj$ (jet cut on the Born) the
channel totals then scattered far beyond their quoted errors, with isolated
spikes in the VEGAS iterations of the channels other than the hard-collinear
one. A per-point dump located the offending points at the cut boundary. There
the real passes the jet cut while every mapped Born fails, or the other way
round, for example $K=0$ with $R$ up to $10^{7}$ at $s_{41}/\hat s\sim10^{-5}$
($4\parallel1$). At these points:
\begin{itemize}
\item the criterion fails and $\rho_\alpha=1/\Nch$;
\item a fraction $(\Nch-1)/\Nch$ of a collinear-singular real goes to channels
      whose mapping does not flatten the $4\parallel1$ singularity;
\item their weights grow like $1/s_{41}$ on a thin region along the cut
      boundary that extends into the collinear limit. The integral is finite
      but the variance is not.
\end{itemize}
Without channels the same real is integrated entirely with the sector
mapping, which flattens it.

The fix computes the fractions and the criterion from the pieces without the
cuts on their mapped Born, as in eq.~\eqref{eq:rho}, while the cuts stay in
$K_\alpha$ in the integrand. Near every limit $R\,W_{ij}\to K^\nc$ whatever the
cuts, so the real is always shared like the counterterm structure that matches
its singularities. In practice the counterterm routine evaluates the pieces
twice: first with \texttt{chan\_nocut} set, with no fills and no scale
variations, giving $K^\nc$ and $K^\nc_\alpha$; then with cuts, giving $K$ and
$K_\alpha$. The fractions are therefore inclusive over the Born-level cuts, for
final- and initial-state radiation alike.

\section{Initial-state partons}
\label{sec:isr}

\begin{enumerate}
\item \textbf{Initial--initial soft dipole.} It has two mappings, $(i,1,2)$
      and $(i,2,1)$ (rescaling of beam 1 or beam 2), each a channel of its
      own. Since the partial-fraction change of 14 November 2026 (separate
      note), the pieces are
      \begin{equation}
        P_a\,B_{12}\big(\kb^{(i,a,b)}\big),
        \qquad P_a = \frac{s_{12}}{s_{ai}\,(s_{1i}+s_{2i})},
      \end{equation}
      so each initial--initial channel carries only the collinear singularity
      along the beam its mapping rescales. Before that change each channel
      carried $\frac12\,s_{12}/(s_{1i}s_{2i})$. In the code the mapping is
      selected by the emitter \texttt{kedip} in \texttt{M2\_S\_g}.
\item \textbf{Hard-collinear channel.} It uses the mapping of the sector,
      $(i_U,i_S,i_B)$, as set by \texttt{initialise\_sector}. For $j$ initial
      and the reference parton $r$ final, this is $(i,r,j)$, not
      $(i,j,r)$. For $j$ initial the hard-collinear channel always merges with
      a soft channel.
\item \textbf{Reference momentum of the azimuth, per channel.} The sector's
      reference can be one of the channel's labels. Each channel therefore
      gets its own, with the rule of the generator: label 1 if no
      initial-state parton is in the mapping, otherwise the smallest
      final-state label not in the mapping.
\end{enumerate}
The forward initial--initial and final--initial mappings already existed for
arbitrary labels; nothing else changes.

\section{Implementation}

\begin{itemize}
\item \textbf{Run card:} \texttt{channel\_decomposition} (default
      \texttt{False}) and \texttt{channel\_delta} (default $1$). The run stops
      if $\delta\le0$ while the flag is on.
\item \texttt{NLO\_Rsub\_template.f}:
      \texttt{setup\_channels} lists the channels (\texttt{chtrip}) and their
      azimuthal references (\texttt{chref}). \texttt{int\_real} applies the
      channel mapping, computes $\rho_\alpha$ from \texttt{KNC} and
      \texttt{KCHNC}, and fills $\rho_\alpha R W$. At the first 20 points it
      checks that the pieces add up to $K$ and writes the result to
      \texttt{testR\_<i>\_<j>.log}.
\item \texttt{NLO\_K\_template.f}: \texttt{local\_counter\_chan} evaluates the
      pieces in two passes and returns $K$, $K_\alpha$, $K^\nc$, $K^\nc_\alpha$.
\item \texttt{NLO\_limits/S\_g.f} and the seven \texttt{HC\_*.f} files:
      dipole and mapping selectors (\texttt{ldip}, \texttt{mdip},
      \texttt{kedip}), the \texttt{chan\_nocut} switch, and no fills at zero
      weight.
\item \texttt{driver\_npo\_template.f}: the channel loop (warm-ups, sharing of
      the calls, sums).
\item \texttt{sectors.py}: the soft and hard-collinear blocks of
      \texttt{local\_counter\_chan}, and \texttt{has\_soft}/\texttt{has\_hc}.
\end{itemize}
With the flag off, the sequence of operations is unchanged: results and
histograms are bit-identical to the code before the change.

The core of \texttt{int\_real} with the flag on:
\begin{lstlisting}
      if(KNC.ne.0d0.and.abs(int_real_no_cnt-KNC).lt.channel_delta*abs(KNC))then
         rho=KCHNC/KNC
      else
         rho=1d0/dble(nchan)
      endif
      ...  fill rho*int_real_no_cnt at the real-emission point
      int_real_orient=(rho*int_real_no_cnt-KCH)*wgt_chan
\end{lstlisting}

\section{Validation}

\paragraph{General checks.}
\begin{itemize}
\item In every sector tested ($e^+e^-\to3,4$ jets, $pp\to e^+e^-$,
      $pp\to Zj$), the pieces add up to $K$ to machine precision
      ($\le2.2\times10^{-16}$).
\item \textbf{Coverage.} Each channel mapping reproduces the integral of a
      smooth test function computed with the sector mapping, within
      $1.5\sigma$:
      \begin{itemize}
      \item $e^+e^-\to4$ jets, sector 3\_6, test function $W$ times the
            Jacobian: 7 channels;
      \item $gu\to Zgu$, test function $RW$ times the luminosity, the
            Jacobian and $\prod s_{ij}/\hat s$;
      \item $u\bar u\to Zgg$, all sectors, test function $W$ times the
            luminosity and the Jacobian (sector 4\_5: $14195\pm31$ with the
            sector mapping, $14244\pm63$, $14191\pm66$, $14200\pm68$,
            $14167\pm59$ with the four channel mappings).
      \end{itemize}
\item With the flag off, all runs are bit-identical to the old code.
\end{itemize}

\paragraph{Final-state radiation, equal calls per sector.}
\begin{table}[h]
\centering
\begin{tabular}{lccc}
\toprule
sector & flag off & flag on & difference \\
\midrule
$e^+e^-\to4j$, 3\_6 (10 seeds) & $-2.70(23)\times10^{-3}$ & $-2.39(9)\times10^{-3}$ & $1.3\sigma$\\
$e^+e^-\to3j$, 3\_4 (6 seeds)  & $-6.9(4.7)\times10^{-4}$ & $-6.0(4.9)\times10^{-4}$ & $0.1\sigma$\\
$e^+e^-\to3j$, 3\_5 (6 seeds)  & $-5.30(25)\times10^{-3}$ & $-5.28(28)\times10^{-3}$ & $0.1\sigma$\\
\bottomrule
\end{tabular}
\caption{Sector totals with the channel decomposition off and on.}
\end{table}

In $e^+e^-\to4$ jets the channels reduce the spread over seeds from
$7.4\times10^{-4}$ (sector 3\_6) and $7.0\times10^{-4}$ (sector 3\_4) to
$2.7\times10^{-4}$ and $1.0\times10^{-4}$. Without channels the spread is
about twice the quoted error, a sign of heavy tails; with channels it matches
the quoted error. In $e^+e^-\to3$ jets the precision is the same. Distribution
means agree ($\chi^2$/bin of on$-$off between 0.4 and 1.6).

With damping factors ($\alpha=2$), $e^+e^-\to jj$ at 1\,TeV gives
$V+R = 0.019915(79)$ with the flag off and $0.019995$ with it on (after the
fix of sect.~\ref{sec:cuts}), against the exact
$\frac{\alpha_s}{\pi}\sigma_B = 0.019977$.

\paragraph{Initial-state radiation.}
\begin{itemize}
\item $pp\to e^+e^-$: the $q\bar q\to e^+e^-g$ sectors have two channels (the
      two initial--initial mappings) and the $qg$ sectors one. Sector means
      agree within $1.3\sigma$; total $R$ is $128.544$ with the flag on against
      $128.526$ with it off. Since no Born cut ever fails here, the uncut
      fractions are bit-identical to the cut ones.
\item $gu\to Zgu$ alone, after the fix of sect.~\ref{sec:cuts}
      ($40\mathrm{k}+80\mathrm{k}$ calls, 2 seeds): pulls on$-$off at most
      $1.9\sigma$; per-channel VEGAS $\chi^2/{\rm dof}$ between 0.5 and 2, with
      no spikes left. Channel errors are 1.2--2 times larger than without
      channels in this subprocess.
\item Full $pp\to Zj$ (82 sectors, user's Mac, seed 0):
      \begin{itemize}
      \item total $R$: $-932.7\pm20.4$ (off) against $-925.5\pm23.9$ (on),
            $+0.2\sigma$;
      \item sector pulls: mean $+0.07$, rms $1.01$;
      \item median ratio of the errors, on/off: $1.00$.
      \end{itemize}
      The largest pulls ($+4.4$, $+3.0$ in $u\bar u\to Zgg$ sectors 4\_5 and
      5\_4) are in final-state sectors with initial-state partons. They have
      the same sign as in the cloud runs, and the coverage test of these
      sectors is complete. We attribute them to the VEGAS bias of low
      statistics per channel (6 sectors $\times$ 4 channels sharing the calls of
      one subprocess); a second seed would settle it.
\end{itemize}

\section{Status and open points}
\begin{itemize}
\item The decomposition is off by default. It is exact by construction, but
      its effect on precision depends on the process: a clear gain in
      $e^+e^-\to4$ jets, neutral in $e^+e^-\to3$ jets and $pp\to Zj$
      overall, and somewhat worse in single initial-state subprocesses.
\item The ISR tests predate the partial-fraction change of the
      initial--initial soft dipole. The pieces still add up to $K$ by
      construction, but the run with the flag on has not been repeated since.
\item Not tested: massive emitters or recoilers in the channel mappings
      ($pp\to t\bar t$).
\item Possible improvements: smoother cuts, to reduce the residual cut-boundary
      variance; a better sharing of the calls among channels.
\end{itemize}

\end{document}
